%
% 3-laurentreihen.tex -- Laurentreihen
%
% (c) 2025 Prof Dr Andreas Müller
%
\section{Laurent-Reihen und meromorphe Funktionen
\label{buch:analytisch:section:laurentreihen}}
\kopfrechts{Laurent-Reihen}
Analytische Funktionen sind immer auf einem Kreisgebiet um den
Entwicklungspunkt definiert.
Der Radius dieses Gebietes wird beschränkt durch die nächstliegende
Singularität der Funktion.
Es kommt jedoch sehr häufig vor, dass eine holomorphe Funktion
nur in einem Ringgebiet definiert ist.
Die einfachsten Beispiele dafür sind die Funktion $z^{-n}$ mit
$n\in\mathbb{N}\setminus\{0\}$, die auf $\mathbb{C}\setminus\{0\}$
definiert sind.
In diesem Abschnitt werden Funktionen auf einem Ringebiet genauer
untersucht und es wird eine Erweiterung der Potenzreihenentwicklung
für negative Potenzen, die Laurent-Reihe, konstruiert.

%
% Funktionen auf einem Ringgebiet
%
\subsection{Funktionen auf einem Ringgebiet}
\input{chapters/040-analytisch/fig/fig-ring.tex}%
Wir betrachten eine holomorphe Funktion $f\colon U\to\mathbb{C}$.
Die offene Menge $U$ sei so gestaltet, dass der Ring
\[
R(z_0;r_1,r_2)
=
\{
z\in\mathbb{C}
\mid
r_1 \le |z-z_0| \le r_2
\}
\]
in $U$ enthalten ist.
Abbildung~\ref{buch:analytisch:laurent:fig:ring} zeigt diese Situation.
Da die Funktion $f$ im Inneren des kleinen Kreises mit Radius $r_1$ um
den Punkt $z_0$ nicht überall definiert sein muss, ist der Integralsatz
von Cauchy nicht auf den grossen Kreis anwendbar.
Der folgende Satz zeigt dass es aber trotzdem möglich ist, den Wert von
$f$ in einem beliebigen Punkt des Ringgebietes mit einem Wegintegral
zu berechnen.

\begin{satz}
Seien $\gamma_2$ und $\gamma_1$ die im Gegenuhrzeigersinn durchlaufenen
Randkreise des Ringgebietes $R(z_0;r_1,r_2)$ mit Radien $r_2$ und $r_1$.
Dann gilt
\begin{equation}
f(a)
=
\frac{1}{2\pi i}
\oint_{\gamma_1}
\frac{f(z)}{z-a}
\,dz
-
\frac{1}{2\pi i}
\oint_{\gamma_2}
\frac{f(z)}{z-a}
\,dz
\label{buch:analytisch:laurent:eqn:zweiintegrale}
\end{equation}
für jeden Punkt $a\in R(z_0;r_1,r_2)$.
\end{satz}

\begin{proof}
Der in Abbildung~\ref{buch:analytisch:laurent:fig:ring} dargestellte
Weg, der vom Punkt $A$ ausgehend zunächst dem äusseren Randkreis
folgt bis wieder zurück zum Punkt $A$, dann radial zum Punkt $a$, in 
einem Halbkreis vom Radius $\varepsilon$ um den Punkt $a$ bis zum
Punkt $B$, von dort im Uhrzeigersinn um den Inneren Randkreis zurück
zum Punkt $B$ und erneut radial zum Punkt $A$, mit einem kleinen, 
kreisförmigen Umweg vom Radius $\varepsilon$ um den Punkt $a$.
Die Funktion ist im Inneren des Weges überall holomorph,
nach dem Integralsatz von Cauchy gilt daher
\[
\oint_{\gamma_1-\gamma_2-\gamma_{\varepsilon}} \frac{f(z)}{z-a}\,dz 
=
0.
\]

Die radialen Teile werden zweimal in entgegengesetzter Richtung durchlaufen
und heben sich daher weg.
Einen Beitrag zum Integral leisten daher nur die Kreise.
Der kleine Kreis um den Punkt $a$ wird mit $t\mapsto \varepsilon e^{2\pi it}$
parametrisiert.
So wird das Integral
\begin{align*}
0
&=
\oint_{\gamma_1} \frac{f(z)}{z-a}\,dz
-
\oint_{\gamma_2} \frac{f(z)}{z-a}\,dz
-
\int_0^1
\frac{f(a+\varepsilon e^{2\pi it})}{a+\varepsilon e^{2\pi it}-a}
\cdot
\frac{d}{dt}
\varepsilon
e^{2\pi it}
\,dt
\\
\oint_{\gamma_1} \frac{f(z)}{z-a}\,dz
-
\oint_{\gamma_2} \frac{f(z)}{z-a}\,dz
&=
\int_0^1
\frac{f(a+\varepsilon e^{2\pi it})}{\varepsilon e^{2\pi it}}
\cdot
2\pi i\varepsilon e^{2\pi it}
\,dt
\\
&=
2\pi i
\int_0^1
f(a+\varepsilon e^{2\pi it})
\,dt.
\end{align*}
Die rechte Seite geht für $\varepsilon\to 0$ gegen $0$, so dass
\[
f(a)
=
\frac{1}{2\pi i}
\oint_{\gamma_1}
\frac{f(z)}{z-a}
\,dz
-
\frac{1}{2\pi i}
\oint_{\gamma_2}
\frac{f(z)}{z-a}
\,dz
\]
folgt.
\end{proof}


%
% Die Laurent-Reihe
%
\subsection{Die Laurent-Reihe}
Auf einem Kreisgebiet lässt sich eine holomorphe Funktion immer in eine
Potenzreihe entwickeln.
Auf einem Ringgebiet ist dies nicht mehr möglich.
An die Stelle der Potenzreihe tritt ein sogenannte Laurent-Reihe.

\begin{definition}[Laurent-Reihe]
\label{buch:analytisch:laurent:def:laurentreihe}
Eine \emph{Laurent-Reihe} ist eine Reihe der Form
\index{Laurent-Reihe}%
\begin{equation}
f(a)
=
\sum_{k=0}^\infty a_k(a-z_0)^k
+
\sum_{k=1}^\infty b_k(a-z_0)^{-k}
\label{buch:analytisch:laurent:def:laurent}
\end{equation}
mit $a_k\in\mathbb{C}$, $k\in\mathbb{N}$ und $b_k\in\mathbb{C}$,
$k\in\mathbb{N}\setminus\{0\}$.
\end{definition}

Die erste Reihe in
\eqref{buch:analytisch:laurent:def:laurent}
ist eine gewöhnliche Potenzreihen, deren Konvergenzradius wir mit $r_1$
bezeichnen.
Die zweite Reihe können wir auch als
\begin{equation}
\sum_{k=1}^\infty b_k(a-z_0)^{-k}
=
\sum_{k=1}^\infty b_k\biggl(\frac{1}{a-z_0}\biggr)^{k}
\label{buch:analytisch:laurent:def:laurent2}
\end{equation}
schreiben, also als Potenzreihe in $1/(a-z_0)$.
Ist $\varrho$ der Konvergenzradius dieser Potenzreihe, dann konvergiert
die Reihe~\eqref{buch:analytisch:laurent:def:laurent2}
offenbar für $|a-z_0| > 1/\varrho = r_2$.
Nach dem Satz von Cauchy-Hadamard ist die Laurent-Reihe ist also genau
dann auf dem Ringgebiet $R(z_0;r_1,r_2)$ definiert, wenn 
\begin{align*}
\frac{1}{r_1}
&=
\limsup_{n\to\infty}\!\sqrt[n]{|a_n|}
&&\text{und}&
r_2
&=
\limsup_{n\to\infty}\!\sqrt[n]{|b_n|}.
\end{align*}
und $r_1 > r_2$.

\begin{satz}[Laurent-Reihe]
\label{buch:analytisch:laurent:satz:laurentreihe}
Ist $f$ eine auf einer Umgebung des Ringgebietes $R(z_0;r_1,r_2)$ definierte
holomorphe Funktion, dann kann sie in eine Laurent-Reihe entwickelt werden.
Ist $\gamma$ ein beliebiger Weg im Ringgebiet, der $z_0$ genau einmal
im Gegenuhrzeigersinn umläuft, dann sind die 
Koeffizienten durch
\begin{align}
a_k
&=
\frac{1}{2\pi i}
\oint_{\gamma} 
\frac{f(z)}{(z-z_0)^{k+1}}
\,dz
&&\text{und}&
b_k
&=
\frac{1}{2\pi i}
\oint_{\gamma} 
f(z)(z-z_0)^{k-1}
\,dz
\label{buch:analytisch:laurent:eqn:laurentkoeffizienten}
\end{align}
gegeben.
\end{satz}

\begin{proof}
Wir berechnen jetzt die beiden Integrale in
\eqref{buch:analytisch:laurent:eqn:zweiintegrale} separat mit
verschiedenen Entwicklungen des Integranden.
Das erste Integral lässt sich wie früher in eine geometrische
Reihe
\[
\frac{1}{2\pi i}
\oint_{\gamma_1}
\frac{f(z)}{z-a}
\,dz
=
\sum_{k=0}^\infty
\biggl(
\frac{1}{2\pi i}
\oint_{\gamma_1}
\frac{f(z)}{(z-z_0)^{k+1}}
\,dz
\biggr)
(a-z_0)^k
=
\sum_{k=0}^\infty a_k(a-z_0)^k
\]
entwickeln.

Für das Integral über den inneren Randkreis 
verwenden wir die geometrische Reihe
\begin{align*}
\frac{1}{z-a\mathstrut}
&=
\frac{1}{z-z_0-(a-z_0)\mathstrut}
\\
&=
-
\frac{1}{a-z_0}\cdot \frac{1}{1-\displaystyle\frac{z-z_0\mathstrut}{a-z_0}}
=
-
\frac{1}{a-z_0}
\sum_{k=0}^\infty \biggl(\frac{z-z_0}{a-z_0}\biggr)^k
\\
&=
-
\sum_{k=0}^\infty \frac{1}{(a-z_0)^{k+1}}(z-z_0)^k.
\end{align*}
Damit kann das Integral über $\gamma_2$ geschrieben werden als
\begin{align*}
-
\frac{1}{2\pi i}
\oint_{\gamma_2}
\frac{f(z)}{z-a}
\,dz
&=
\frac{1}{2\pi i}
\oint_{\gamma_2}
\sum_{k=0}^\infty
f(z)(z-z_0)^k
\frac{1}{(a-z_0)^{k+1}}
\,dz
\\
&=
\sum_{k=0}^\infty
\biggl(
\frac{1}{2\pi i}
\oint_{\gamma_2}
f(z) (z-z_0)^k
\,dz
\biggr)
\cdot
\frac{1}{(a-z_0)^{k+1}}.
\intertext{Durch Verschiebung des Summationsindex wird daraus}
&=
\sum_{k=1}^\infty
\biggl(
\frac{1}{2\pi i}
\oint_{\gamma_2}
f(z) (z-z_0)^{k-1}
\,dz
\biggr)
\cdot
\frac{1}{(a-z_0)^k}
=
\sum_{k=1}^\infty
b_k
\cdot
\frac{1}{(a-z_0)^k}.
\end{align*}

Die beiden Reihenentwicklungen können wir jetzt zusammensetzen zu
\begin{align*}
f(a)
&=
\sum_{k=0}^\infty
a_k
\cdot 
(a-z_0)^k
+
\sum_{k=1}^\infty
b_k
\cdot 
(a-z_0)^{-k},
\end{align*}
die versprochene Laurent-Reihe.

Jetzt muss noch gezeigt werden, dass die Integrale über $\gamma_1$ 
bzw.~$\gamma_2$ durch Integrale über den Weg $\gamma$ irgendwo
im Ringgebiet ersetzt werden können.
Die Integranden der Integrale sind im ganzen Ringgebiet holomorph.
Die Deformation des Weges $\gamma_1$ in den Weg $\gamma$ oder des
Weges $\gamma_2$ in den Weg $\gamma$ ändert daher das Integral nicht.
\end{proof}

%
% Singularitäten
%
\subsection{Singuläre Punkte}
Die Laurent-Reihenentwicklung ermöglicht, eine Funktion in der Umgebung
eines singulären Punktes genauer zu untersuchen.
Sei $U$ ein offene Teilmenge von $\mathbb{C}$ und $z_0$ ein isolierter
Punkt von $\mathbb{C}\setminus U$.
Ein auf $U$ definierte, holomorphe Funktion $f\colon U\to\mathbb{C}$
ist also in jeder genügend kleinen Umgebung von $z_0$ in jedem Punkt
definiert ausser im Punkt $a$.
Eine solche Stelle heisst ein \emph{isolierter} singulärer Punkt.
\index{singularer Punkt@singulärer Punkt}%

Die Funktion $f$ ist in jedem Ringgebiet $R(z_0;r_1,r_2)$, um den Punkt
$z_0$ definiert, sofern nur $r_1$ klein genug ist und für 
alle inneren Radien $r_2<r_1$.
Nach Satz~\ref{buch:analytisch:laurent:satz:laurentreihe} gibt es
ein Laurent-Reihenentwicklung
\begin{equation}
f(z)
=
\sum_{k=0}^\infty a_k(z-z_0)^k
+
\sum_{k=1}^\infty b_k(z-z_0)^{-k},
\label{buch:analytisch:laurent:eqn:laurentsing}
\end{equation}
die für $z$ beliebig nahe an $z_0$ konvergiert.
Es folgt, dass die Reihe
\[
u(z)
=
\sum_{k=1}^\infty b_kz^k
\]
für beliebig grosse $z$ konvergiert, sie ist also in ganz $\mathbb{C}$
definiert.
Am Nullpunkt ist $u(0)=0$.
Die zweite Summe in der Laurent-Reihe
\eqref{buch:analytisch:laurent:eqn:laurentsing}
kann daher als $u(1/(z-z_0))$ geschrieben werden.
Die Funktion $u(1/(z-z_0))$ heisst der \emph{singäre Teil} der Funktion $f$
in einer Umgebung von $z_0$.
\index{singularer Teil@singulärer Teil}%

Falls $u=0$ ist, dann handelt es sich um eine \emph{hebbare} Singularität,
die Funktion kann stetig in den Punkt $z_0$ hinein fortgesetzt werden.

\begin{definition}[Ordnung eines singlären Punktes]
Sei $z_0$ ein singulärer Punkt der Funktion $f(z)$ und $u(z)$ der 
singuläre Teil an der Stelle $z_0$.
Wenn $u(z)$ ein Polynom vom Grad $n$ ist, dann heisst $z_0$ ein
\emph{Pol der Ordnung $n$}.
\index{Pol}%
Falls $u(z)$ kein Polynom ist, heisst $z_0$ eine \emph{essentielle
Singularität}.
\index{essentielle Singularitat@essentielle Singularität}%
Falls $u=0$ ist und $a_k=0$ für $k\le n$ in 
\eqref{buch:analytisch:laurent:eqn:laurentsing},
dann heisst $z_0$ eine \emph{Nullstelle der Ornung $n$}.
Die Funktion $\omega(f,z_0)$ ist definiert durch
\[
\omega(z_0;f)
=
\begin{cases}
-\infty&\qquad\text{falls $z_0$ eine essentielle Singularität von $f$ ist}\\
-n&\qquad\text{falls $z_0$ ein Pol der Ordnung $n$ von $f$ ist}\\
\phantom{-}m&\qquad\text{falls $z_0$ eine Nullstelle der Ordnung $m$ von $f$ ist}\\
\phantom{-}\infty&\qquad\text{falls $f=0$ ist}
\end{cases}
\]
\end{definition}

Für die Ordnung $\omega(z_0,f)$ gelten die folgenden Rechenregeln.
\begin{align*}
\omega(z_0,f+g) &\ge \min\bigl(\omega(z_0,f),  \omega(z_0,g)\bigr)
\\
\omega(z_0,fg) &= \omega(z_0,f) + \omega(z_0,g).
\end{align*}

Wenn $f$ eine Singularität der Ordnung $n$ an der Stelle $z_0$ hat, 
dann ist $(z-z_0)^{-n}f(z)$ eine analytische Funktion mit einer
hebbaren Singularität an der Stelle $z_0$, die dort nicht verschwindet.
Insbesondere ist dann auch $1/(z-z_0)^{-n}\cdot 1/f(z)$ eine in einer
Umgebung von $z_0$ holomorphe Funktion, so dass man
\[
\omega(z_0;f)
=
-\omega(z_0;1/f)
\]
schliessen kann.

%
% Meromorphe Funktionen
%
\subsection{Meromorphe Funktionen}
Funktionen mit lauter isolierte Singulariäteten sind besonders
leicht zu behandeln und verdienen daher einen eigenen Namen.

\begin{definition}[meromorphe Funktion]
Eine holomorphe Funktion $f\colon U\to\mathbb{C}$,
für die $\mathbb{C}\setminus U$ aus lauter isolierten Punkten besteht,
heisst \emph{meromorph}.
\index{meromorph}%
\end{definition}

Die wichtigsten Beispiel für meromorphe Funktionen sind die rationalen
Funktionen
\[
f(z)
=
\frac{f(z)}{g(z)},
\]
wobei $f(z)$ und $g(z)$ Polynome sind.
Die Singularitäten sind dann genau die Nullstellen von $g(z)$ und die
Ordnung der Pole von $f(z)$ ist die Ordnung der Nullstellen von $g(z)$.


